\section{Mechanism of Forking}
\label{sec:mechanism}

The approximate power law in Figure~\ref{fig:frequency} is our first clue: why should context frequency organize the gap so simply?

\subsection{From the frequency law to novel continuations}
\label{sec:coverage}
A rare context has had few chances to reveal its possible continuations. Call a validation target \emph{novel} when its context occurs in training but was never followed there by that target, and let $f(F)$ be the fraction of validation positions at context count $F\ge1$ whose continuation is novel. This fraction closely tracks the frequency dependence of the gap: the dashed curves in Figure~\ref{fig:frequency} are the directly measured $f$, rescaled and shifted,
\begin{equation}
\gap_t(F)\approx a_t f(F)+b_t,
\label{eq:coverage-main}
\end{equation}
with one scale and one offset per branch at training step $t$. Good--Turing estimation connects this novel fraction to singleton counts in the corpus; heavy-tailed conditional distributions provide a reason for its slow, approximately power-law decay. The detailed derivation and the terms beyond simple proportionality are in Appendix~\ref{app:math}. Thus the frequency trend points to a property of the text: rare contexts leave more plausible continuations unseen. It remains to explain why those continuations are harmed by replay.

\subsection{Crowding-out under replay}
\label{sec:sharpening}
A finite corpus does not fully represent the continuation distribution of each context: a context whose population distribution assigns equal probability to $C$ and $D$ may appear three times with $C$ and once with $D$ in training, so fitting the sample underpredicts $D$ on fresh data. The replay trajectory adds a further clue. After one complete pass, the corpus and its novel fraction are fixed, yet the gap continues to grow, by comparable amounts over the first few passes (Figure~\ref{fig:staircase}; Appendix~\ref{app:replay-clue}). The missing coverage alone does not specify this continuing damage. We hypothesize that each repetition reinforces the seen continuations and progressively suppresses the alternatives.

This would make forking ordinary overfitting confined to very small problems. For a context $c_i$, call the next-token logits computed from it its \emph{context subspace}, and split the vocabulary into the continuations of $c_i$ seen in training, those that appear only in validation, and the rest (Appendix~\ref{app:crowding-derivation}). If training shows a single continuation $C$, the loss in that subspace is $-\ln p_C$; an idealized step directly in logit coordinates raises the logit of $C$ by $\eta(1-p_C)$ and, through the softmax, lowers every other logit by $\eta p_v$. Mass flows from the unseen continuations onto the seen one, so the seen-to-unseen odds grow by a factor above one at every pass with nothing to stop them while $p_C<1$ (Appendix Equation~\eqref{eq:zt}). We call this \emph{crowding-out}; it provides an idealized account of the growth in the novel-continuation penalty with replay.

\begin{figure}[!ht]
\centering
\includegraphics[width=\linewidth]{figures/fig_replay_transfer_schematic.pdf}
\caption{Local sharpening and cross-context recovery, illustrated schematically. Dashed marks show prior heights; darker extensions show gains. \emph{Others} denotes the mean probability of the top 20 remaining tokens.}
\label{fig:replay-transfer}
\end{figure}

Figure~\ref{fig:replay-transfer} illustrates the proposed mechanism: replay sharpens a local distribution, while updates from related contexts can support alternative continuations. Its heights are illustrative, and the no-gram-only arrow idealizes the contrast rather than asserting zero coupling in $n$-gram models. Appendix~\ref{app:sharpening} presents the measured context distributions, including ``walk their dog,'' together with aggregate seen/novel losses and predictive entropy. The aggregate measurements show worsening novel-continuation accuracy in the both-table model while the no-gram control improves; individual contexts need not sharpen monotonically. Crowding-out also offers an explanation for continued forking after the tables are frozen (Section~\ref{sec:interventions}): the context-specific representation remains available, and backbone updates can continue to favor the observed continuations.

This leaves one question: why does the same replay improve the no-gram model on novel continuations, while the $n$-gram model grows less accurate? The next experiment tests whether the table allows sharpening to proceed more independently within each context.

\FloatBarrier
\subsection{Testing the locality of context subspaces}
\label{sec:context-locality}
The mechanism of Section~\ref{sec:sharpening} suggests that the $n$-gram module gives each context a coordinate on which sharpening can accumulate. How independently do these coordinates evolve? Context-loss masking tests whether removing a context's own training loss suppresses further gap growth while other contexts keep training. Appendix~\ref{app:context-locality} presents the experiment, its controls, and the measured continuation masses and gap contributions.

The original trajectories and paired one-pass analyses suggest that direct replay supervision contributes to local gap accumulation: removing it can suppress average growth on selected contexts while growth on unmasked peers continues. Individual masked contexts nevertheless change, and other masking runs have not recovered the same aggregate behavior under different measurement protocols (Appendix~\ref{app:mask-reconciliation}). These observations support investigating context-specific sharpening, while independence across contexts remains unestablished.

\FloatBarrier
